\documentclass{ltjbook}
\usepackage{docmute} 
\usepackage{../myStyle} 

\begin{document} 
\chapter{古典的テスト理論} 
\label{chapter6} 

この講義の目的である心理統計は，\textbf{「見えないもの」をみよう・測ろう・考えようとする}学問です。見えないものを見るためには，測定に対する考え方，方法論の整備が必要です。ここでは測定に対する考え方と基本的なモデルについて解説します。
\section{テストと心理学}
心理学の研究対象である心は，見えないし触れることもできません。なんなら存在するかどうかもわかりません。我々は自分も含めて，他者，人間の中には「心がある」と思っていますが，物質的なものとは違って本当に存在するのかどうかを確認する方法がありませんから\footnote{目に見えている物体だってそこにあるかどうかわかりません。カントはその著書「純粋理性批判」\citep{カント1961jo,カント1961ge}の中でものが見えて，存在するとはどういうことかを考え続けました。興味がある人は入門書として\citet{黒崎政男2000}などを読んでみるといいでしょう。}，いわば「心がある」と考えることにしているにすぎません。逆に言えば，それで存在することで合意するなら，ペットにも心があると言っていいでしょう\footnote{ペットは家族の一員である，などと言いますが，それはその動物にも自我があると家族全員がみなしてふるまっているのですから，少なくとも社会的実在としてその家庭の中ではペットに自我があることになります。もっとも，動物嫌いの人にとっては他人の愛玩動物であっても獣の一種にしか見えないでしょうし，法律的にも愛玩動物を傷つけることは障害罪ではなく器物破損罪として扱われる事実があります。}。ともかくここでのポイントは，心の存在は「みなし」あるいは「仮定」にすぎない，ということです。心とは，機能的な面を表現した\keyterm{モデル}{もでる}{model}なでしょう。

モデルとは，模型とか理想形という意味でもありますが，ここでは抽象的に考えられた機構(しくみ)のことを指します。人の心は，人を理解するために我々が想定したモデルだというわけです。モデルはそれがどのように振る舞うかを記述したものですが，記述の方法が数学的であれば\keyterm{数理モデル}{すうりもでる}{}と言いますし，確率を取り入れたモデルであれば\keyterm{確率モデル}{かくりつもでる}{}，統計的な指標同士の関係を記述したものであれば\keyterm{統計モデル}{とうけいもでる}{}などと言います。記述は自然言語でも構いませんが，第\ref{chapter1}講でもお話しした通り，関係を数式で表現する\keyterm{数学的現象主義}{すうがくてきげんしょうしゅぎ}{}は近代科学の基本的なスタイルです。数学的に表現することで，厳密さや抽象化，一般化が進みます。

さて心理学の研究パラダイム，あるいは心理学の基本的な心のモデルとは何かといえば，\pageref{fig::01_01}の図\ref{fig::01_01}に示した，\keyterm{SR図式}{SRずしき}{}です。刺激と反応の組み合わせから，ブラックボックスである人の中に「心」というモデルがあることを考えるわけです。この刺激と反応のペアについて，心理学は様々な方法を試してきた学問であるといえるでしょう。一方で，同じような枠組みで見えないものを見ようとする試みがあります。みなさんにも馴染み深い，「テスト」と呼ばれるものです。

みなさんはこれまで，様々なシーンでテストを受けてきたと思います。そのほとんどは，学校関係での学力テストでしょう。
ところで改めて考えると，「学力」とはなんでしょうか。文部科学省は学習指導要領という教育の指針のなかで，様々な力を育てることを謳っていますが，「学ぶ力」とか「生きる力」といったものは明確に定義できているようで，よくわかりません。
ただ少なくとも，学力テストでは学力を測っているということだけは間違いありません。学力は目に見えない，個人のうちに秘めたる力であり，これを測定するためのテストは心理学とも深い関係を持っています。

学力テストは，紙とペンを使って行うのがほとんどです。テスト問題用紙には，様々な設問が用意されており，受験生はそれに回答していきます。回答は正答と誤答に区分され，正答数が多いと学力が高く，正答数が低いと学力が低いと判断されます\footnote{実際には，正答数を数えるだけではなく，問題の重要度などに応じて配点を変えて合計することがあります。}
つまり，同じ人に対する様々な設問という刺激を与え，正答という反応がたくさん帰ってくると，その人の中に設問を正しく捉えて判断するための「知識」「学力」「能力」という一貫した性質が保持されていると考えているのです。反応のパターンが同じであるということは，ブラックボックスにその性質が宿っていると考えるのは，心理学のモデルそのものですね。

\subsection{人の性格についての考え方}
例えば心理学では，人の性格(人格)というものを「個人の持つ安定的な行動傾向」と捉えたりします。この「安定的ななにか」がある，と考えるところが重要で，安定した性質であれば，複数の刺激に対して同じパターンの回答を返すだろう，ということが予想できます。性格検査はそのような目的で作られています。

性格心理学は，古くは外見などの見た目で人を判断するところから始まりました。筋肉質な体つきだから真面目だろうとか，太っているから朗らかだろう，という思いつきのようなところから始まっています。ガルというひとが「骨相学」を考えましたが，これは頭蓋骨の形を見て性格を判断しよう，というものです。これは偏見や差別にもつながる話であり，今では完全に否定された考え方ですが，歴史的にはそうした時代があったのです\footnote{詳しくは\citep{渡辺茂2019}を参照。}。ところで，こうした分類で考える方法を類型論といいます。類型論は話が単純になっていいのですが，このカテゴリーに収まらない例外もすぐに思い付いてしまいます。そこで性格心理学では，類型論から特性論へとシフトしてきました。特性論とは，性格のような人を特徴付ける性質は人類全体で共通していて，そんの程度が個々人によって違うだけだ，と考えるのです。例えば明るい人，暗い人というのがいますが，これは「外交性」という性質が高い人なのか，低い人なのかという違いであって，暗い人は外交性スコアの低い人，明るい人は外交性スコアの高い人，という相対的な関係で理解するのです\citet{小塩真司2020-08-25}。

このような特性を測定するのが，性格検査と呼ばれるものです。これはテスト理論の性格版ともいえます。学力テストと違うのは，測りたい目に見えないものが一種類ではないということです。例えば国語のテストでは国語の能力を測ります。算数のテストでは算数の能力を測ります。つまり測りたいものが1次元なのです。これに対して性格は，様々な行動に反映される個人の特性ですから，それが何次元あるのかが問題になってきます。性格検査はその歴史と統計的な分析の結果から，性格は5次元で測るのがよかろう，とひとまず結論づけられています。ここでは学力テストと性格テストが同じ目的を持ったツールであること，ただし主眼点がことなることを踏まえておいてください。

\subsection{良いテストと悪いテスト}
ところで，学校で行われる試験の中には，学力の程度を明確に測定することを目的にしたものもありますが，授業内容や基礎的な知識を有しているかどうかといった理解度・達成度を目的にしているものもあります。小学校の低学年などはえてして後者の目的で作られており，満点が出ることも少なくありません。ここでは前者の，学力の高さ低さを精緻に測定することを目的にしたテストについて考えていきます\footnote{いわゆる心理テストと銘打って，「うさぎさんを選んだからあなたは心の底では助平なのです」といったお遊びとは違うことに注意してください。そうしたお遊びは，箔をつけるために心理学という学問の名前を利用しているにすぎず，その分類や予言?にいかなるデータも根拠もありませし，ましてや何かを測定することを目的にしたものではありません。そのような遊びをする人はよく，これは統計だから，と言ったりしますが，科学的な手続きが取られたものでもなければ，データとして蓄積されたものが数量的に評価できず，統計的に分析できるわけでもありません。いちいちこのような断りを入れなければならないのは，個人的には甚だ迷惑で，こちらの営業妨害だと言いたくなります。やれやれ。}。

さて，目に見えない能力を性格に測定することを目的にしたテストを考えた場合，理想はどういうテストになるでしょうか。考えやすくするために，理想の逆，最悪のテストとは何か，ということを考えてみましょう。みなさんが今まで受けてきた試験で，これは良くないテストだな，と思ったことはありませんか？それはどのようなものでしたか？
回答例はいくつか挙げられると思います。1つは「不公平なテスト」でしょう。同じ回答をしていても，ある人は正答と判断され，ある人は誤答と判断される，というのでは困ります。これは問題と採点者の両方が原因です。まず問題に対して，正答が厳格に定義されていることが必要です。そして採点者は，回答者が誰であってもその「正答の定義」に則って正しく判断する必要があります。記述・論述試験は，正答の定義や採点者の裁量に影響されることがありますから，本来あまり望ましいテストの形式ではありません。

他にも，「問題が間違っていて答えがない」ものも良いテストとは言えないですね。こうしたテストは全員誤答になるか全員正答になるしかありません。その場合，テストの点数は受験者全員で一律になり，\keyterm{分散}{}{variance}が$0$になります。分散が$0$というのは，そこから情報を得ることができないので，よくないテストだということになります。

逆に言えば，良いテストとは評価基準が一律で，誰がどこで誰に対して採点しても正答・誤答の判断が一位に定まり，かつ，そのテストから得られる分散がなるべく大きいものであること，と言えます。もしテストが一問しか作れないのであれば，受験生の半数が正答し，残る半数が誤答するのが良いテストです。それが分散を最大にするからです。みなさんはAkinator\footnote{\url{https://jp.akinator.com/}}というのをご存知でしょうか。古くからある20Quesutions(20の扉)というゲームの，オンライン版なのですが，これはいくつかの設問に対する反応から答えを当てるというゲームです。Akinatorは，我々が考えたことに次々と質問してきます。それに答えていくと，20問と行かないうちに当てられることがほとんどです。このAkinatorの出題方法をみていればわかりますが，その問いは「出題者の考えそうな問題空間」を二分の一に分割していくようなものを選んでいます。例えば第一問は「女性ですか？」と聞かれたりします。これにYesと答えるかNoとこたえるかで，少なくとも半分に絞り込めますね\footnote{もちろん性別は2種類に限りませんし，Akinatorへの回答選択肢は「はい」「いいえ」「わからない」「多分そう/部分的にそう」「多分違う/そうでもない」になっていますが，ここでは話を簡単にするために回答をYes/Noの二択に限定しています。}。次に「名前に漢字が入っている？」など，さらに回答を絞り込むようにしています。この絞り込み方が1/2のペースでいけば，一番効率よく世界を狭められるのです\footnote{Akinatorを実際にやってみると，二問目に「Youtuber？」とか「闘う？」とか聞かれたりします。これはAkinatorの持っている回答データベースがこれに偏っているか，あるいはエンターテイメント性を出すためにあえて1/2に拘っていないかといった可能性が考えられます。}。人間がこの手のゲームをやると，下手な質問としていきなり答えを特定するような質問をしたりします。例えば2,3問目に「あなたの想定した答えは``綾瀬はるか''ですか？」と聞いて，YesならいいのですがNoであれば何も分かったことになりません。綾瀬はるかでない何かのほうが，選択肢が多いからです。良いテストとはこのように，世界を半分にわけ，それをさらに半分にわけ・・・と段階的に区分できるような問題が次々あるもの，ということになります。もし1回の試験で様々な区分を作りたいのであれば，事前に分かっている問題の難易度順に1,2,3...と並べ，正答の程度が回答者に応じて分布するようにすれば大きな分散＝情報を得ることができます。

このように，テストは目に見えないものを測定するためのモデルであり，それを数理的に記述して特徴を明らかにしようという営みがあります。これを\keyterm{テスト理論}{てすとりろん}{Test Theory}といいます。次に\keyterm{古典的テスト理論}{こてんてきてすとりろん}{Classical Test Theory}について解説していきます。

\section{測定のモデル}
古典的テスト理論は，テスト理論のスタートとなった(今とは古典である)理論です。そのモデルは非常に単純な形で表現できます。古典的テスト理論の測定モデルは次のようになります。

\[ X = t + e \]

なんとシンプルなのでしょう。ここで$X$とはテストの点数を表します。また$t$は\keyterm{真のスコア}{しんのすこあ}{true score}，$e$は\keyterm{誤差}{ごさ}{error}と言われます。$t$は本当に測りたいもの，例えば本当の学力，真の実力といったもの，と思ってください。古典的テスト理論のモデルが言っていることは，テストの点数は本当の能力を反映してはいるが，それに誤差がついたものになっているよ，ということです。この誤差都心のスコアの分離，が重要なポイントです。たとえばハードサイエンス(物理学など)では，測定は心を持たないものを対象にし，また測定精度の高いセンサー等で測定しますので，測定した値は大体その通りで間違い無いと考えます。体重計が83kgだったら体重は83kgであって，60kgや90kgということはありえそうにありません。ですが，学力や性格のようなものは，心を持った(と想定される)ものの反応です。人間は嘘をつくし，間違えるし，易きに流れる生き物ですから，外交的なふりをして答えようと振る舞うこともできますし，うっかり計算ミスをしてしまうこともありますし，「あなたはだらしないところがあるのではないか」といわれたら「そうかもなあ」とおもって反応してしまうかもしれません。そういう意味で，表に出てきた回答（反応やテストの得点）はそのまま信用できないのです。本当は$t$なんだけど，なんらかの事情である$e$が加わったものが，目に見える$X$なんだよ，と考えるのが，このモデルの意味するところです。

測定されたものをそのまま信用しない，といえば聞こえが悪いですが，測定の精度はそれほど高くないのだということを自省的に取り込んでいるという意味で，ソフトサイエンスの良心があらわれているとも言えます。みなさんは人を対象に研究をするのですから，この点を強く自覚しておいてください\footnote{個人的な経験ですが，他分野の人，とくに理系領域から来られた人と話をする時に，測定されたものを盲信的に使おうとする態度に驚かされることがあります。心理学者としては当然「回答者が答えたからと言ってその通りとは限らない」と考えるのですが，回答がそのまま正しいと思う人もいます。それらの数字は間違っているかもしれず，回答者の心理的な要因によって数字がハックされることもある，ということは常に注意した方がいいでしょう。数値目標を建てると，裏技的な攻略方法が編み出されるため，測定が意味をなさなくなることがある，ということを\citep{ジェリー2019}は詳しく論じています。}。


\section{誤差の基本的性質}
ここで誤差のことを考えてみましょう。誤差は本当に測定したいものからずれている分，を表します。
例えば体重80kgの人がいたとして，その人が体重計に乗ったら81kgだったとします。このずれた1kgがこの測定の誤差になります。
ところで，この体重計に誰も乗っていない時に，その目盛りが0.5kgを表示していたらどうでしょうか。つまり，何もなくてもプラス0.5kgになるようにずれているわけですから，誰が乗ってもプラス0.5kgのずれが常に生じることになります。これは体重計が壊れているので，ゼロにリセットできるように調整が必要です。このような，常に生じる誤差のことを特に\keyterm{系統誤差}{けいとうごさ}{systematic error}といいます。このような誤差が紛れ込むのはある程度仕方がないものかもしれませんが，常にずれているのであれば，その分を差し引くなどして補正することができる誤差です。

ではこの人の，残りの0.5kgのずれはなんでしょうか。例えば測定の前にたくさん食べたとか，たまたま厚手の下着を身につけていたとかいった，その場限りの一貫性のない誤差というのもあり得ましょう。こう言った誤差のことを\keyterm{偶然誤差}{ぐうぜんごさ}{random error}と言います。この誤差は，どうしてそのようになったのか，その理由がわかりません。「たまたま」ですから，たまたまプラスになる要因もあれば，たまたまマイナスになる要因もあるでしょう。この偶然誤差こそ，測定の時に入り込む困った問題です。系列誤差は調整して直すことができますが，偶然誤差は傾向がないので対応のしようがないからです。

テストの時に考えなければならない誤差$e$は，基本的にこの偶然誤差です。たまたま前日ヤマを張ったところが出たとか，たまたま欠落した単語帳のページに書いてあった項目が問題に入ったとか，体調がたまたま良かった・悪かった，と言った要因など，真の実力が発揮できない要素は必ず入ってくると考えなければなりません。

さてでは，テスト理論のモデルに戻ってこのことを考えてみましょう。もとのモデルを少し拡張して，個人$i$についてのテストであるということを明確にするために，$X_i$という添字をつけて書くとします。
\[ X_i = t_i + e_i \]
さて，このテストを$n$人のクラスで実施したとします。そこからクラスの平均点が算出できますね。
\[ \bar{X} = \frac{1}{n}\sum_{i=1}^n x_i \]

これを展開すると次のようになります。
\begin{align*}
    &       \bar{X}  =  \frac{1}{n}\sum_{i=1}^n X_i  & &\mbox{}\\
    &                    =  \frac{1}{n}\sum_{i=1}^n (t_i + e_i) &\mbox{定義より}\\
    &                     =  \frac{1}{n} \sum_{i=1}^n  t_i + \sum_{i=1}^n  e_i &\mbox{分配して}\\
    &                     =  \bar{t} + \bar{e} &\mbox{}\\
\end{align*}

さて，ここでポイントです。$\bar{e}$は偶然誤差の平均を表しています。偶然誤差は一貫性のない誤差なので，プラスに出るかマイナスに出るかわかりませんが，\textbf{理屈上，その平均は$0$である}と考えられます。これがなんらかの定数$c$であれば，それは系列誤差であり補正すべき大きさですが，そうではなく全く何の傾向もないのですから。この特徴，$\bar{e}=0$を代入すると上の式は，

\[ \bar{X} = \bar{t} \]

となります。やったね！すごい！誤差が消えて無くなりました。

このように，測定の平均値は真の値の情報だけになることがわかります。測定すれば真の値にたどり着く，というのが古典的テスト理論のポイントなのです。

ちょっと補足をしておきます。
最終的に得られた$\bar{X} = \bar{t}$というのは，テストの平均が真の値の平均になるということですが，それはクラス全体の平均が本当の能力だというだけで，個人の能力が本当かどうかはわからないじゃないか，と思われるかもしれません。それはそうなのですが，この話を個人と学級という対比ではなく，テストの複数の項目と，合計点としてのその人のスコアだと考えてみてください。すなわち，$X_{ij}$として，個人$i$の項目$j$に対する正答，誤答が採点され，最終的に$X_i = \sum_{j=1}^M X_{ij}$と点数が与えられるとします(項目数は$M$個)。

ここでも真の能力$t_i$と，誤差が考えられますが，誤差は$e_{ij}$と項目ごとに逐一入り込むと考えられます。すると上の式展開は次のようになります。
\begin{align*}
&           \bar{X_i}  =  \frac{1}{M}\sum_{j=1}^M X_{ij} \\
&                         =  \frac{1}{M}\sum_{j=1}^M (t_i + e_{ij}) \\
&                         =  \frac{1}{M}\sum_{j=1}^M t_i + \frac{1}{M}\sum_{j=1}^M e_{ij} \\
&                         =  t_i + \bar{e_{i}} \\
\end{align*}
ここで$\bar{e_i}$は個人に生じる誤差の平均ですが，これも特に傾向が考えられないので$0$と考えられますから，$\bar{X}_i = t_i $と個人の真の能力を測っている，と考えることができるでしょう。



\section{テストの分散と信頼性・妥当性の関係}
\subsection{信頼性}
\subsection{妥当性}



% \item[誤差の基本的性質]測定に誤差が含まれている可能性については，ティコ・ブラーエ，ケプラーなどによって，天文学における膨大なデータと，それに対する天体運動の数式の当てはめを考えていた頃から懸念されていた問題であった。ガウスは誤差のありうるべき形を考え，そこから正規分布を導出しており，これが古典的テスト理論はもちろん，心理統計学における誤差や心理的実態についてのモデルに援用されることになる。ここでは誤差の分布として平均が0の正規分布について理解する。
% \begin{flushright}
%     $\to$ \citet{長沼201611}のP.16--40，あるいは少し数学的な説明が長くなるが，\citet{西内201712}のP.408--430
% \end{flushright}

% \item[テストの平均値と真値の関係]古典的テスト理論における測定に含まれる誤差が，各項目について同様のモデルで考えられるとすると，平均値の定理および誤差分布の性質から，複数の項目を測定することで真値が得られると考えられる理論的根拠が導出される。ここではこの過程を確認するとともに，類似した項目を幾度となく繰り返して測定する尺度がなぜそのような形式を取ることになっているのか，その原理について理解する。

% \item[テストの分散と信頼性・妥当性の関係]古典的テスト理論のモデルを分散について同様の式展開を行う。ここでは誤差の新たな性質である，他の何者とも相関関係にないことから，テストで得られる分散が真の分散と誤差の分散「だけ」に分割されることを理解し，テストにおける信頼性の定義を行う。また目に見えないものを測定しようという営みである限り，信頼性だけでなく妥当性についても考えなければならないが，数値化できる信頼性とは違って妥当性は論証するしかないこと，信頼性は妥当性の上限であることを理解し，心理学における測定の重要性を再確認する。
% \begin{flushright}
%     $\to$ テストの信頼性については上述の\citet{kosugi2016b}のP.20--25。
% \end{flushright}



\end{document} 
